\section{The Makropoulous Affair}

Janacek's\footnote{\url{https://www.youtube.com/watch?v=nmHoYKjEXCs}} famous opera\footnote{\url{https://en.wikipedia.org/wiki/The\_Makropulos\_Affair\_\%28opera\%29}} centers around the discovery of a secret elixir which can actually grant the brewer immortality. It is an ancient Greek recipe from the dawn of time:

\begin{quotex}
During the end of the 16th century Emperor Rudolf II\footnote{\url{https://en.wikipedia.org/wiki/Emperor\_Rudolf\_II}} gave his alchemist Hieronymus Makropulos the task of preparing a potion that would extend his life and ordered him to test it on his daughter Elina. She fell into a coma and Hieronymus was sent to prison. But after a week Elina woke up and fled with the formula, starting a new, itinerant life becoming one of the best singers of all time.
\end{quotex}


After a trial settles a contest over a will between two intertwined families, one of the actors is revealed as having multiple identities over centuries. When their identity is remembered and revealed by chance, ``Elina'' gives up the recipe, and dissolves rapidly into old age and death before their eyes as she recites the Lord's Prayer. Those who behold this, burn the recipe.

Bernard Williams\footnote{\url{https://en.wikipedia.org/wiki/Bernard\_Williams}}, a brilliant and erudite ``psuedo-conservative'' (a member of that innumerable but perishing set of geniuses who instinctively resisted the worst consequences of Revolution, but effectively tolerated the advance guard) wrote a famous essay\footnote{\url{http://www.unc.edu/~jfr/RI-TMCR1.htm}} on this opera, a meditation on death and mortality. He concluded that physical immortality was an undesirable choice for anyone who wished to overcome the tragedy of the human condition.

It is instructive to watch this ``half-way house'' analysis of the human condition -- negative examples can teach a lot, especially when they are allowed to so dramatically run rampant (as in our day, the end of days). Aside from the Christian overtones of the ending, a traditionalist might suggest that the point is being missed by both the author and the analyst, and indeed by everyone but perhaps the composer (music being our last link to the divine in a material age, since the ``hearing is the last thing to go'').

Yes, physical immortality, without regeneration and perfection, is actually a curse -- as both Methodius of Tyre\footnote{\url{https://en.wikipedia.org/wiki/Methodius\_of\_Olympus}} and Paul Althaus\footnote{\url{https://en.wikipedia.org/wiki/Paul\_Althaus}} concur in observing, God's imposition of death upon fallen man was actually an act of mercy. To have lived forever as we began would be not immortality, but eternal death. Given this, our age's obsession with eternal youth, comfort, and insulation is merely a psychotic avoidance. Thus, alienation proceeds from insulation, which proceeds from an even profounder, hidden alienation, that of the self from the self. This is the mystery that Christianity has begun to forget or to obscure.

Since, certainly, to burn the recipe is to preserve humankind from a great deal of harm in the short run (much like taking Vaclav Havel's Forgetting We Are not God\footnote{\url{http://www.orthodoxytoday.org/articles5/HavelGod.php}} speech as one's spiritual \emph{magna carta}), many are eager to try this. But this is to deny the possibility of redemption as well, for the removal of harm results in the further privation of man.

Even worse, this gesture is futile, since in the cosmic drama, the recipe which is burned naturally has a second copy -- perhaps someone memorized it -- or an angel re-dictates it in a dream? Thereby the physical copy can no more perish than the immortal article, for the two are in a sense one. In defiance of this, historic Christianity today is forever preferring either the afterlife, or this one. This is un-Christian.

God swears a double oath in Hebrews\footnote{\url{https://books.google.com/books?id=BHl-a3\_Z1OYC\&pg=PA154\&lpg=PA154\&dq=hebrews+double+oath\&source=bl\&ots=O91NPd\_lmA\&sig=wdMm0vSsRghvkmk-LyyIEFeFLrw\&hl=en\&ei=15p6ToreH8qksQK-uvTdAw\&sa=X\&oi=book\_result\&ct=result\&resnum=1\&ved=0CB4Q6AEwAA\#v=onepage\&q\&f=false}} (significantly mentioning Melchizedek), on earth as it is in heaven, just as in the Lord's Prayer. This is the macro-microcosm, and it renders the earth a theater for conflict, for the taking of heaven by storm. In this conflict, there are heroes and there are traitors, champions and rabble, sinners and saints. Trying to render life ``safe'' is a disservice -- burning the recipe for immortality ought only properly be done by those who realize that they are not thereby destroying it.

Thus Hermeticism attempts to reconcile the pagan emphasis on this world with the Christian accent upon the future goal of the afterlife. The two were meant to be one, if perhaps secretly, and by using depth, rather than surface dialectic, we see that physical immortality and future eternal life may be tragically opposed for now, yet remain inseparable. As Wilhelm von Humboldt\footnote{\url{https://www.amazon.com/Humanist-Without-Portfolio-Anthology-Writings/dp/B000IY152K}} might say, ``each is too fine to sacrifice to the other in totality''. This is what a well-bred man from past centuries would have said, when confronted with choosing between Christianity and paganism.

A Hermeticist might suggest that the burning of the recipe was incense offered to heaven, and will result in the denial of this secret to the lustful, and the proffering of it to the worthy, through miraculous and yet physically articulate means (after all, God understands His own miracles, so they are not inherently un-rational). Rather than the cycle of endless repetition or the opposing angles of dialectic, he will see convergence of spirals which emerge from immutable patterns in the dark and tend towards an epiphany beyond imagining.

Without this depth, there is no redemption, only tragedy. There is no hope. There is no Graal\footnote{\url{https://www.youtube.com/watch?v=fjB2n\_9809c}}. And there can be no summons to take heaven by storm, because we will either take sides with greedy physical immortalists or pious after-lifers, and ignore what is right before us.

Gornahoor is right to summon the ancestral spirits, including the Christians, to its side. Like Aragorn at the end of \emph{The Lord of the Rings}, we will have to tread the path of the dead\footnote{\url{https://en.wikipedia.org/wiki/Paths\_of\_the\_Dead}}, and hold ourselves to a new, unbreakable oath. That which has been separated will be re-united. The sword which was broken will be re-forged.


\flrightit{Posted on 2011-09-23 by Logres}

\begin{center}* * *\end{center}

\begin{footnotesize}
\begin{sffamily}
\texttt{Boreas on 2011-09-25 at 16:38 said:}

Although this comes from completely different direction, smoehow this reminds me ofa science article I read the other day while visiting the local library. An american scientist named Raymond Kurtzweil has proposed that man's personal consciousness and memory can be imprinted into electric chips by the year 2049 and then man will have practically gained ``immortality''. He has many dazed followers in the US who listen to him like some sort of a prophet. He himself eats 200 vitamin etc. pills a day and follows a very strict dietary regimen to be able to live to that day, so he can then ``live forever''.

This is complete, total madness, magic of the blackest kind if succesful. I feel very sorry for the people who believe the stuff this false prophet ``showing people grest signs and wonders'' offers as a panacea for human mortality. Tenth level of Dante's Inferno would be a paradise compared to living in an electric chip with a personal consciousness and memory for aeons. It is truly an utterly sad and desperate age when men start clinging on their petty lives like this.

\hfill

\texttt{logres on 2011-09-25 at 17:07 said:}

Every traditionalist I am around (Christian or not) has the exact, precise same reaction to Kurzweil; otherwise very intelligent and sophisticated people go in for this kind of stuff in a mind-blowingly naive way. It makes the alchemical quest look like solid science. An acquaintance of mine via the Net swears by this guy, \& thinks ``singularity'' is almost upon us. Which, if it is, suicide becomes a moral option, because this would be far worse than the occult.

\hfill

\texttt{Boreas on 2011-09-27 at 05:34 said:}

You have some fine `traditionalists' around you, logres. This kind of men and theories should be criticized vehemently and without mercy with a ``latter day intellectual militancy'', as Charles Upton has called it at Sophia Perennis.

The fact that ``otherwise inteligent and sophisticated people'' fall into this kind of mania is a very telling sign of the fact that wordly intelligence and outward sophistication are nothing in themselves, and can even be a serious hindrance in the spiritual path. ``It is very difficult for a rich man (in this case, for a one who is content with one's self and knowledge, and who is not ready to leave `life' behind) to enter into the kingdom of heaven''.


\hfill

\end{sffamily}
\end{footnotesize}

